% ==============================================================================
% TEMPLATE: BAB HASIL DAN PEMBAHASAN
% ==============================================================================
% File ini memuat struktur bab Hasil dan Pembahasan.
% WAJIB terbagi menjadi \section{Hasil Praktikum} dan \section{Tugas Praktikum}
% dengan pemisah \clearpage tepat di antaranya.
% Dokumentasi gambar wajib memuat 'Penjelasan Gambar :' dengan poin Input & Output.
% ==============================================================================

\chapter{Hasil dan Pembahasan}
\label{chap:hasilpembahasan}

% ==========================================
% HASIL PRAKTIKUM (PANDUAN TERBIMBING)
% ==========================================
\section{Hasil Praktikum}
\label{sec:hasilpraktikum}

Pada bagian Hasil Praktikum, dilakukan serangkaian tahapan praktikum terbimbing sesuai dengan panduan modul [Nama Topik Praktikum]. Kegiatan ini mencakup persiapan lingkungan kerja, pemuatan data atau dependensi, implementasi algoritma inti, hingga analisis terhadap luaran sistem.

\subsection{Inisialisasi Lingkungan dan Persiapan Modul}
Tahap awal dimulai dengan mempersiapkan seluruh modul, berkas dependensi, serta konfigurasi perangkat lunak yang dibutuhkan selama proses praktikum.

\begin{figure}[H]
    \centering
    \includegraphics[width=0.55\textwidth]{1.png}
    \caption{Source Code Inisialisasi dan Import Library}
    \label{fig:1}
\end{figure}
Penjelasan Gambar :
\begin{itemize}
    \item \textbf{Input:} Mengeksekusi perintah impor pustaka standar manipulasi data dan visualisasi grafik yang relevan dengan topik praktikum.
    \item \textbf{Output:} Seluruh modul berhasil diinisialisasi ke dalam memori lingkungan komputasi tanpa terjadi konflik dependensi atau kesalahan jalur pustaka.
\end{itemize}

\subsection{Implementasi dan Prapemrosesan Data}
Tahapan inti dilakukan untuk mentransformasikan data masukan serta mengeksekusi logika pemrosesan utama.

\begin{figure}[H]
    \centering
    \includegraphics[width=0.50\textwidth]{2a.png}
    \caption{Source Code Prapemrosesan Data}
    \label{fig:2a}
    
    \vspace{0.2cm}
    
    \includegraphics[width=0.50\textwidth]{2b.png}
    \caption{Output Verifikasi Struktur Data Hasil Transformasi}
    \label{fig:2b}
\end{figure}
Penjelasan Gambar :
\begin{itemize}
    \item \textbf{Input:} Pada Gambar \ref{fig:2a}, baris kode memproses variabel masukan melalui tahapan normalisasi dan restrukturisasi format agar sesuai dengan standar masukan modul.
    \item \textbf{Output:} Gambar \ref{fig:2b} memvalidasi dimensi larik data dan ringkasan nilai statistik yang menunjukkan bahwa data telah siap diproses pada tahapan pemodelan/eksekusi selanjutnya.
\end{itemize}

\subsection{Visualisasi dan Evaluasi Luaran}
Kinerja luaran dianalisis secara visual dan kuantitatif guna memverifikasi ketepatan hasil eksekusi program.

\begin{figure}[H]
    \centering
    \begin{minipage}{0.48\textwidth}
        \centering
        \includegraphics[width=\textwidth]{3a.png}
        \caption{Visualisasi Luaran Metode 1}
        \label{fig:3a}
    \end{minipage}\hfill
    \begin{minipage}{0.48\textwidth}
        \centering
        \includegraphics[width=\textwidth]{3b.png}
        \caption{Visualisasi Luaran Metode 2}
        \label{fig:3b}
    \end{minipage}
\end{figure}
Penjelasan Gambar :
\begin{itemize}
    \item \textbf{Output:} Gambar \ref{fig:3a} dan Gambar \ref{fig:3b} memperlihatkan komparasi visual hasil pengujian secara berdampingan. Terlihat bahwa implementasi Metode 2 menghasilkan stabilitas kurva yang lebih adaptif terhadap pola data dibandingkan Metode 1.
\end{itemize}

% ==========================================
% TUGAS PRAKTIKUM (STUDI KASUS MANDIRI)
% ==========================================
\clearpage
\section{Tugas Praktikum}
\label{sec:tugaspraktikum}

Pada bagian Tugas Praktikum, dilakukan penyelesaian studi kasus mandiri menggunakan dataset atau skenario baru yang dirancang untuk menguji pemahaman mahasiswa terhadap konsep yang telah dipelajari.

\subsection{Deskripsi Penugasan dan Konfigurasi Parameter}
Studi kasus ini menuntut analisis mandiri terhadap skenario [Jelaskan skenario penugasan mandiri secara ringkas].

\begin{figure}[H]
    \centering
    \includegraphics[width=0.55\textwidth]{T1.png}
    \caption{Source Code Implementasi Solusi Tugas Praktikum}
    \label{fig:t1}
\end{figure}
Penjelasan Gambar :
\begin{itemize}
    \item \textbf{Input:} Menerapkan algoritma solusi pada dataset tugas mandiri dengan konfigurasi hyperparameter yang telah dioptimalkan.
    \item \textbf{Output:} Menampilkan log eksekusi komputasi yang mengonfirmasi bahwa seluruh tahapan penugasan telah terselesaikan dengan sukses.
\end{itemize}

\subsection{Hasil Pengujian dan Analisis Performa}
Evaluasi komprehensif dilakukan untuk menguji sejauh mana solusi yang diusulkan mampu menyelesaikan permasalahan yang diberikan.

\begin{figure}[H]
    \centering
    \includegraphics[width=0.60\textwidth]{T2.png}
    \caption{Output Hasil Pengujian dan Metrik Evaluasi Tugas Praktikum}
    \label{fig:t2}
\end{figure}
Penjelasan Gambar :
\begin{itemize}
    \item \textbf{Input:} Perintah evaluasi metrik kinerja pengujian terhadap himpunan data uji mandiri.
    \item \textbf{Output:} Tabel atau log nilai metrik performa yang menunjukkan bahwa solusi mandiri berhasil mencapai tingkat keandalan yang memenuhi kriteria kelulusan pengujian modul.
\end{itemize}

\subsection{Analisis Komparatif dan Pembahasan}
Berdasarkan serangkaian percobaan yang telah dilakukan pada studi kasus tugas mandiri, diperoleh beberapa wawasan penting:
\begin{enumerate}
    \item \textbf{Kesesuaian Pendekatan Terhadap Karakteristik Masalah:} [Analisis teknis mengapa metode yang dipilih efektif]...
    \item \textbf{Efisiensi Waktu Komputasi dan Penggunaan Sumber Daya:} [Evaluasi kecepatan eksekusi dan trade-off kompleksitas]...
    \item \textbf{Potensi Pengembangan Lanjutan:} [Peluang perbaikan atau optimasi parameter lebih lanjut di masa depan]...
\end{enumerate}
